\section{Определение билинейной формы}

\begin{definition}[Билинейная форма]
Пусть $V$ --- линейное пространство над $\mathbb{R}$ ($F = \mathbb{R}$).

Функция $f(x, y): V \times V \to \mathbb{R}$ называется \textbf{билинейной формой}, если выполнены следующие 4 свойства:
\begin{enumerate}
    \item $f(x_1 + x_2, y) = f(x_1, y) + f(x_2, y),\quad \forall x_1, x_2, y \in V$;
    \item $f(\lambda x, y) = \lambda f(x, y),\quad \forall x, y \in V,\ \forall \lambda \in \mathbb{R}$;
    \item $f(x, y_1 + y_2) = f(x, y_1) + f(x, y_2),\quad \forall x, y_1, y_2 \in V$;
    \item $f(x, \mu y) = \mu f(x, y),\quad \forall x, y \in V,\ \forall \mu \in \mathbb{R}$.
\end{enumerate}

В координатах:
\[
f(x, y) = \sum_{i,j=1}^{n} a_{ij} \xi_i \eta_j,
\]
где $\xi_i, \eta_j$ --- координаты векторов $x, y$ в базисе $E$.
\end{definition}
